\section{Etude du code FreeRTOS}

\subsection{Clignotement de la LED (Partie 3.4)}

À partir du code fourni en annexe et en suivant les étapes de création du projet sur Microchip Studio, nous avons observé le bon clignotement de la \textit{LED0}. \\

La tâche \textit{Horloge} définit un délai d'attente de 1000 ms (\textit{HORLOGE\_TASK\_DELAY 1000}). L'état de la LED changeant toutes les secondes, un cycle complet (allumé puis éteint) dure 2 secondes. Le système fonctionne donc correctement avec une fréquence de clignotement confirmée de 0,5 Hz.

\begin{lstlisting}[style=code]
void vCodeBlinkLED (void * pvParameters)
{
    /* declaration de variables locales de la tache */
    bool Blink;
    
    /* initialisation des variables locales */
    Blink = false;
    
    for (;;)
    {
        /* attente du semaphore d'eveil */
        xSemaphoreTake (Sem, portMAX_DELAY);
        
        /* Cycle H*/
        if (Blink == false) {
            /* turn LED on. */
            port_pin_set_output_level (LED_0_PIN, LED_0_ACTIVE);
            Blink = true;
        } else {
            /* turn LED off. */
            port_pin_set_output_level (LED_0_PIN, LED_0_ACTIVE);
            Blink = false;
        }
    }
}
\end{lstlisting}

Le résultat de cette exécution sur la carte SAMD21, confirmant la fréquence de clignotement, est visible via ce lien : \href{https://youtu.be/XMjjLOAeuDM}{\textbf{voir la démonstration vidéo sur YouTube}}.

\subsection{Réponses à l'étude du code FreeRTOS (Partie 3.5)}

\begin{itemize}
    \item \textbf{Procédure d'ordonnancement :} Dans le code générique (\textit{tasks.c}), la procédure chargée du changement de contexte et de la sélection de la prochaine tâche à exécuter est \textit{vTaskSwitchContext()} (qui utilise la macro \textit{taskSELECT\_HIGHEST\_PRIORITY\_TASK()}). \\
    
    \item \textbf{Lien avec \textit{FreeRTOSConfig.h} :} La politique d'ordonnancement est directement pilotée par les constantes de ce fichier, principalement \textit{configUSE\_PREEMPTION} (qui active/désactive le mode préemptif) et \textit{configMAX\_PRIORITIES}. \\
    
    \item \textbf{Adaptation de la politique :} Adapter fondamentalement la politique (par exemple vers un ordonnancement de type EDF) n'est pas trivial. Il faudrait directement modifier le code source du noyau dans \textit{tasks.c}, en particulier la procédure \textit{vTaskSwitchContext()} et l'organisation des listes de tâches (\textit{pxReadyTasksLists}). \\
    
    \item \textbf{Politiques de gestion de la mémoire (heap\_1 à heap\_5) :}
    \begin{itemize}
        \item \textit{heap\_1.c} : Allocation uniquement, aucune libération possible (déterministe, idéal pour les systèmes statiques).
        \item \textit{heap\_2.c} : Allocation et libération autorisées, algorithme \textit{best-fit}, mais sans fusion des blocs libres (risque de fragmentation).
        \item \textit{heap\_3.c} : Encapsule les fonctions standards \textit{malloc()} et \textit{free()} du compilateur pour les rendre \textit{thread-safe}.
        \item \textit{heap\_4.c} : Allocation et libération avec fusion automatique des blocs adjacents libres (défragmentation).
        \item \textit{heap\_5.c} : Similaire à \textit{heap\_4.c}, mais permet de répartir le tas (heap) sur plusieurs zones mémoires non contiguës.
    \end{itemize}
\end{itemize}
